\section{Introduction}
\label{sec:intro}

A production-deployed LLM agent rarely fails by forgetting an anchor outright. It
fails by a more subtle failure mode: the model continues to \emph{verbally}
mention an anchor --- ``warehouse at $(12,8)$'', ``defend territory~$X$'', ``the
patient is on warfarin'' --- in its running summaries and reflections, while
its \emph{action-emitting} directives quietly stop reflecting it. The summary
still says ``I should defend territory~$X$''; the next dispatched unit goes
elsewhere. The reflection still references the warehouse; the next pick task
allocates zero weight to deliveries there. We call this failure mode
\textbf{stale coherence}.

Stale coherence is the failure mode that breaks long-running agent deployments,
and yet it is invisible to the dominant evaluation regimes. Existing memory
benchmarks measure only one axis of the phenomenon. We argue this is a single,
operationally definable phenomenon whose joint behaviour distinguishes it from
every individual axis those benchmarks measure.

\paragraph{Existing benchmarks measure only one axis.}
\begin{itemize}[leftmargin=*]
\item \textbf{Verbal-recall benchmarks.}
LongMemEval~\citep{Wu2025LongMemEval}, RULER~\citep{Hsieh2024RULER},
Lost-in-the-Middle~\citep{Liu2024LostInMiddle}, and
InfiniteBench~\citep{Zhang2024InfiniteBench} measure whether the model can
\emph{retrieve} a fact from long context, not whether it sustains
\emph{action} on it. Their probe is verbal: ``what was X?''
\item \textbf{Action-grounded benchmarks (session-discrete).}
MemoryArena~\citep{He2026MemoryArena} and
MemoryAgentBench~\citep{Hu2025MemoryAgentBench} do score action grounding,
but in session-discrete agent loops: each task is bounded, the agent has
explicit reset boundaries, and the action set is small. They cannot expose a
drift that emerges only across continuous tick-by-tick context accumulation.
\item \textbf{Production agent observations.}
FAIR Cicero~\citep{FAIR2022Cicero} and Anthropic's multi-agent research
system~\citep{AnthropicMultiAgent2025} anecdotally describe stale coherence
(``the agent kept saying it would defend territory~$X$ but stopped moving
units there''; inter-channel coherence loss in long-running sessions), but
provide no controlled measurement.
\end{itemize}
No existing benchmark dissociates verbal mention from sustained behavioural
commitment on tick-continuous, naturally-accumulated context. The phenomenon
falls in the gap between three literatures and consequently has neither an
operational definition nor a measurement instrument.

\paragraph{Contributions.}
We make three contributions, organized phenomenon-first:
\begin{description}[leftmargin=*]
\item[C1 --- Phenomenon.] An operational definition of stale coherence as the
dual probe
$\bigl(\,\text{verbal\_recall}(t),\,\text{action\_grounded\_recall}(t)\,\bigr)$
with dissociation metric
$\Delta(t) = \text{verbal\_recall}(t) - \text{action\_grounded\_recall}(t)$.
Stale coherence is distinct from forgetting (both axes can remain non-zero),
distinct from positional bias (the anchor is fixed at $t = 0$ and probed at
the same surface continuously), and distinct from session-discrete action
grounding (the agent never receives a query about the anchor; behavioural
commitment is measured implicitly across emission cadence).

\item[C2 --- Method.] Anchor injection at $t = 0$ with a paired (verbal token,
implicit action goal) --- e.g.\ (``warehouse at $(12,8)$'',
\texttt{\{deliver,\ warehouse\}}); a multi-cadence structured-output
substrate (four parallel decision channels emitting at three nested cadences
over a fixed four-tool surface) as the elicitation harness; and 8-sim-hour
episodes matching the natural TTL spread of agent directives (8--24\,h),
which the route-B 240-sim-second framing did not exercise. The substrate is
load-bearing methodologically because single-cadence benchmarks cannot
dissociate verbal mention from sustained behavioural commitment.

\item[C3 --- Empirical (camera-ready).] Cross-family decay profiles of
$\Delta(t)$ over 8 sim-hours, and a statistical-independence test of $\Delta$
against LongMemEval-style verbal-only recall, across at least three LLM
families (Claude-Sonnet-4.6, GPT-5-mini, Llama-3.1-8B).
\end{description}

\paragraph{Single pre-registered hypothesis.}
We register one testable hypothesis driving the entire empirical contribution:
\textbf{H}: for every LLM family $m$ in our model menu, there exists
$t^{\star}(m) \leq 4\,\text{sim-hours}$ such that
$\Delta(t^{\star}) > 0.3$. We additionally test whether
$\Delta$ is statistically independent of LongMemEval-style verbal recall:
Spearman correlation between the dissociation gap and verbal-only recall
$< 0.5$. A correlation above this threshold would reduce stale coherence to
known long-context bias; a correlation below it establishes stale coherence as
a separate axis production agent designers must explicitly test for.

\paragraph{Why stale coherence dissociates from existing memory benchmarks.}
LongMemEval and RULER measure single-axis verbal recall; MemoryArena and
MemoryAgentBench measure action grounding but in session-discrete tasks with
explicit query-response boundaries. Stale coherence is the \emph{joint}
behaviour of both axes under tick-continuous long-horizon context, and it can
appear even when each axis alone looks healthy:
$\text{verbal\_recall}(t)$ near 1 (the model still talks about the anchor) and
$\text{action\_grounded\_recall}(t)$ near 0 (no directive weights it). The
theoretical underpinning for the decay shape itself is the diminishing-returns
geometry described by \citet{Sinha2025IllusionDiminishingReturns}: recall
scores decay sub-linearly with effective context length; stale coherence is
the operational counterpart for the \emph{action} axis and is observable only
when the two axes are measured jointly on naturally-accumulated context.

\paragraph{Roadmap.}
Section~\ref{sec:related-work} surveys memory benchmarks along the two axes
above and the production-agent failure literature;
Section~\ref{sec:architecture} describes the multi-cadence structured-output
substrate that makes the dual probe measurable;
Section~\ref{sec:metrics} presents the dual probe methodology and supporting
scoring stack;
Section~\ref{sec:experiments} reports the cross-family decay experiments;
Section~\ref{sec:discussion} discusses implications for production agent
design;
Section~\ref{sec:limitations} states limitations;
Section~\ref{sec:reproducibility} describes the reproducibility package.
